\section{Conclusion}

In this chapter, we have proposed to extend the framework of constant delay in action execution.
More precisely, we have considered the possibility that an action's effect may spread over several future steps.
Leveraging previous literature on \gls{mc} of higher order, we have proposed adapting the \gls{mtd} and \gls{mtdg} models to \glspl{mdp} in order to obtain \glspl{mtdmdp}.
This choice was made for the great modelling capacity of \glspl{mtd} and their low parameter dimensionality. 
This new model of higher-order \glspl{mdp} offers exactly the possibility of observing the effect of an action over several future timesteps. 
Four variations of the process have been presented, depending on whether it is based upon \gls{mtd} or \gls{mtdg} and depending on the state in which the actions are applied. 
The processes are \glspl{ismmdp}, \glspl{immmdp}, \glspl{psmmdp} and \glspl{pmmmdp}.

We then conducted a theoretical analysis of their properties.
First, we have shown that the \glspl{mtdmdp} can be cast back to \glspl{mdp} by augmentation of the state, as in the constantly delayed case.
We have then focused on learning optimal policies for these models. 
Considering \glspl{psmmdp} and \glspl{pmmmdp}, we showed that an undelayed policy yields the same average reward in these processes as in the underlying undelayed \gls{mdp}.
Then, we have studied the more interesting model of \gls{ismmdp} which contains constant delays as a particular case.
First, the peculiarities of the model have been studied, highlighting how the choice of a delay vector $\boldsymbol\lambda$ acts on the average reward and state distribution of the agent.
A conjecture has been made on the effect of $\boldsymbol\lambda$ on the aforementioned criterion with some supporting arguments.
Then, the effect of the delay on the variance of these quantities has been studied for constantly delayed \gls{dmdp}. 
Finally, we have studied how the properties of the underlying \gls{mdp} might be passed to the \gls{ismmdp}.
These results are of interest to rule out algorithms from the theoretical \gls{rl} literature that cannot be readily applied to \gls{ismmdp}. 
For an algorithm that can be applied to them, UCRL2, we have provided experiments with varying $\boldsymbol\lambda$ to show its effects on the average reward.
In the same idea, we provided another experiment with expected discounted return as an objective to highlight the effect of $\boldsymbol\lambda$ in this case.

Finally, \gls{immmdp} have been left aside for future works.
Considering other future directions, one possibility would be to improve the theoretical \gls{rl} algorithms applied to \gls{mtdmdp} by leveraging the peculiarities of the estimation of the parameters of an \gls{mtd} model \cite{berchtold2002mixture}.
Moreover, our setting is reminiscent of \cite{romano2022multi}, which considers temporally-partitioned rewards, which could be seen as a form of multiple reward delays. 
Studying the connections between these settings might yield better guarantees for theoretical approaches.